%% kapa: exam=1 round=31 year=2020 period=2 form=A subject=accounting subject_ko=회계학 seq=32 number=72
%% source: Q-Net 감정평가사 기출문제 (cst00309) artlSeq=5207219 file=extracted/1008/5207219/A형/A형/2020년도 제31회 감정평가사 1차 시험_2교시_A형.pdf
%% conversion: kordoc 4.24.51 -> tools/gen1.py (auto); review=proofread-sonnet; flags=none
\begin{question}{72}
(주)감평은 두 개의 제조부문(P1, P2)과 두 개의 보조부문(S1, S2)을 두고 있다. 각 부문간의 용역수수관계는 다음과 같다.
\par{\small\centering\begin{tabular}{|c|c|c|c|c|}\hline
\multirow{2}{*}{\makecell{사용부문\\제공부문}} & \multicolumn{2}{c|}{보조부문} & \multicolumn{2}{c|}{제조부문} \\ \cline{2-5}
 & S1 & S2 & P1 & P2 \\ \hline
S1 & - & 50\% & 20\% & ? \\
S2 & 20\% & - & ? & ? \\ \hline
부문발생원가 & ￦270,000 & ￦450,000 & ￦250,000 & ￦280,000 \\ \hline
\end{tabular}\par}
(주)감평은 보조부문의 원가를 상호배분법으로 배분하고 있다. 보조부문의 원가를 배분한 후의 제조부문 P1의 총원가가 ￦590,000이라면, 보조부문 S2가 제조부문 P1에 제공한 용역제공비율은?
\begin{choices}
  \item 20\%
  \item 25\%
  \item 30\%
  \item 35\%
  \item 40\%
\end{choices}
\end{question}
\answer{5}
